% 中文摘要（与 main.tex 的英文论文配套）
% 编译：xelatex abstract_zh.tex
\documentclass[11pt,a4paper]{ctexart}
\usepackage[margin=2.5cm]{geometry}
\usepackage{amsmath}
\usepackage[hidelinks]{hyperref}

\title{\bfseries 基于多客户端 vfio-user 后端的\\同主机虚拟机热迁移中 USB 存储设备透明直通}
\author{匿名投稿\\ \small 稿件编号：\texttt{usbvfiod-usb-migration}\\ \small 按双盲评审要求隐去作者姓名、单位与联系方式}
\date{}

\begin{document}
\maketitle

\section*{摘要}

将物理 USB 设备直通给虚拟机在存储、加密狗与仪器等场景中很有价值，但它与热迁移存在根本冲突：主机侧的设备会话（打开的文件描述符、已声明的接口、端点状态以及设备自身的内部状态）无法序列化，因此任何会拆除该会话的迁移路径都会迫使客户机（Guest）重新枚举设备。本文以 Cloud Hypervisor 与 \texttt{usbvfiod}（一个以 Rust 实现、模拟 xHCI 控制器的 vfio-user 服务端）为具体研究对象。

我们首先通过实测指出：迁移一台挂载了 vfio-user 设备的虚拟机时会发生\textbf{死锁}而非明确报错——目标端 VMM 在协议版本握手中阻塞了 199.9 秒，其间无任何日志、无错误输出。根因是 vfio-user 服务端仅调用一次 \texttt{accept()}，而该连接仍被源端持有，目标端的连接虽然在内核层面建立成功（监听队列 \texttt{Recv-Q}=1），却永远不会被服务。我们还通过源码分析说明：即使解决了连接问题，该设备的 \texttt{Pausable}/\texttt{Migratable} 均为空实现、其快照不含 BAR 与设备内部状态，因此设备状态本身也不会被迁移。

在此基础上，本文提出一套\textbf{无需修改 VMM} 的设计。其关键观察是：vfio-user 的握手类命令（\texttt{Version}、\texttt{DeviceGetInfo}、\texttt{DeviceGetRegionInfo}、\texttt{GetIrqInfo}）根本不访问设备后端，只有数据面命令才访问。因此只要把后端锁从“按连接持有”改为“\textbf{按命令加锁}”，第二个客户端即可在第一个仍然连接的情况下完成握手，而设备状态始终留在共享后端中。围绕这一改动，我们还发现并修复了交接过程暴露出的两个进一步缺陷：（1）即将退出的源端会下发破坏性命令，覆盖存活端刚注册的中断线并解除其 DMA 映射；（2）在中断线完成切换之前的窗口内完成的传输，其完成中断被打到即将退出的 VMM 的事件描述符上而永久丢失，表现为复制静默卡死。我们分别以“中断归属者”规则与“注册新中断线后补发一次中断”加以修复。

测量方法学上，我们记录了若干必须修正的观测陷阱：目标端会以 \texttt{File::create()} 截断串口日志、迁移会重建串口设备从而重置 Guest 的 TTY、以及只读且跳过日志回放的挂载会隐藏最近的写入。因此本文的结论一律取自\textbf{Guest 自身的日志}，而非宿主侧观测；并强调时序相关的缺陷必须通过重复实验而非单次通过来验证。

在真实 32\,GB exfat U 盘直通给 Ubuntu 22.04 客户机的实验环境中，一次跨越迁移点的 128\,MiB 文件复制在 \textbf{20/20} 次运行中成功完成：实测停机 \textbf{4--21\,ms}（中位数 6\,ms；目标 $\le$2000\,ms），数据校验和与源文件\textbf{逐字节一致}，Guest 侧\textbf{零重枚举、零复位、零 I/O 错误}。为把"修复起作用"与"这几次恰好没事"区分开，我们补齐了五类对照：\textbf{不迁移对照}（8/8 通过，复制耗时与迁移臂同量级；无迁移可判，故仅做校验和判定）、\textbf{release 构建臂}（8/8 通过，复制快约 4.5 倍）、\textbf{关闭踢中断的负对照}（7/8 通过，1 次复制停在固定字节数且此后再无完成事件）、\textbf{朴素热拔插基线}（3/3 失败：Guest 内核记录 \texttt{USB disconnect} 与 I/O 错误，\texttt{dd} 返回非零、校验和不符），以及\textbf{故障注入套件}（用两个仅 debug 可见、运行时可开的环境钩子人为放大交接窗口）：\textbf{关闭归属守卫 0/5 通过}（Fisher 精确检验 $p<0.0001$，Guest 出现 36\,s 的 \texttt{DID\_TIME\_OUT} 与重复 I/O 错误，这一项是确定性的）；500\,ms 窗口下开启踢中断 5/5 通过、关闭 4/10 通过（$p=0.044$，但在 6 个对比的 Holm 校正下不显著，仅作名义参考）；\textbf{把窗口放大到 5\,s 后，开启踢中断 8/8 通过、关闭 2/8 通过}（$p=0.007$，对 5\% 与 80\% 的效应量该样本的检验效力为 0.88）——踢中断的必要性由受控放大窗口确立，但即使 5\,s 窗口也并非必然失败，这与“丢失的中断只有当它恰为最后一个未完成命令时才致命”一致。

本文同时如实报告了尚未消除的限制：迁移会重置 Guest 串口控制台（属于 VMM 串口设备模型行为，与 USB 通路无关）；尚未实现显式的传输静默期（quiesce），因此丢失的中断是否致命取决于它是否恰为最后一个未完成命令（自然窗口是否非空取决于暂停时刻的确定方式：把 CH 自己的时钟与 usbvfiod 墙钟对齐后为 20 次中 4 次、4 个完成事件；用 harness 的迁移 epoch 锚定为 8 次、10 个；用迁移请求锚定则是 11 次、24 个的严格上界。三者在同一量级）；跨主机迁移需要协议级的设备状态区域，当前协议不携带所需状态。

\section*{关键词}
虚拟化；热迁移；设备直通；vfio-user；USB；xHCI；Cloud Hypervisor；KVM

\vspace{1em}
\noindent\rule{\textwidth}{0.4pt}
\vspace{0.5em}

\noindent\textbf{英文摘要见} \texttt{main.pdf}（\texttt{paper/main.tex}，IEEE 会议格式）。

\end{document}
